% q4.tex — 七、问题四的分析（2026-08-15, 四段式论证, 素材见 PAPER_HANDOFF.md §8）

问题四要求给出商超还应采集的数据及其对补货与定价决策的帮助。本文的论证按四段式
展开：先指出现有模型的简化假设，再说明所需数据的采集方式，接着给出将数据引入
模型后的目标函数或约束的修正形式，最后分析修正效果。六类建议数据的要点见表10，
其中两类给出详细的公式修正。

\begin{table}[H]
\centering
\caption{建议采集的数据及其对模型的修正}
\label{tab:q4data}
\small
\begin{tabular}{p{2.6cm}p{4.2cm}p{4.6cm}}
\toprule
数据 & 针对的模型简化 & 引入模型后的修正 \\
\midrule
每日库存与废弃量 & 假设剩余商品全部打折售出 & 引入浪费因子 $L$，目标函数扣除废弃损失 \\
单品保质期 & 未区分短保与耐储品 & 引入保鲜系数 $BX$，收紧短保品补货上限 \\
天气与客流 & 需求仅由星期效应解释 & 需求预测乘以弹性修正因子 \\
竞争对手定价 & 加成率 $\eta_c$ 取历史常数 & 加成率随对手价格动态调整 \\
供货商数据 & 假设任意数量即时到货 & 增加起订量、提前期与准时率约束 \\
关联食材销售 & 定价只考虑自身需求 & 需求曲线引入品类间相关系数修正 \\
\bottomrule
\end{tabular}
\end{table}

\subsection{每日库存与废弃量}

现有目标函数式\eqref{eq:obj}假设损耗部分全部按折扣 $d_c$ 清仓回收。实际经营中，
部分蔬菜因品相问题直接废弃，既无收入还需承担清运成本。商超可结合 POS 库存盘点
与废弃称重记录采集每日废弃量，定义浪费因子 $L_c$ 为废弃量占陈列量的比例，则
式\eqref{eq:obj}修正为

\begin{equation}
\max\ Z=\sum_{t=1}^{7}\sum_{c=1}^{6}x_{c,t}\left[P_{c,t}\bigl((1-\bar\ell_c-L_c)+\bar\ell_c d_c\bigr)-\hat w_{c,t}-L_c c_d\right],
\label{eq:q4_waste}
\end{equation}

其中 $c_d$ 为单位废弃处置成本。加入 $L_c$ 后，补货量对预期销量 $B$ 不再"多多益善"，
模型会自动在缺货损失与废弃损失之间取平衡，补货方案更贴近真实经营约束。

\subsection{单品保质期}

当前模型对同一品类内的单品未区分保鲜能力。采集各单品的保质期数据（供应商标签
或系统字段），定义保鲜系数 $BX_i\in(0,1)$（越接近 1 越耐储）。对短保品，其
隔日不可售的比例更高，补货上限应收紧为

\begin{equation}
x_i\le BX_i\,\hat B_i/(1-\ell_i),
\label{eq:q4_bx}
\end{equation}

对耐储品则允许在批发价低位时超出当日需求补货，形成跨日库存。修正后短保品的
浪费风险下降，耐储品的采购成本可随价格波动优化。

\subsection{天气与客流}

现有销量预测式\eqref{eq:weekday_forecast}只利用了星期效应。蔬菜需求对天气与客流
敏感（雨天客流下降、高温带动叶菜消费），可通过气象接口与门店客流计数采集历史
数据，回归得到弹性因子，将需求预测修正为

\begin{equation}
\hat B'_{c,t}=\hat B_{c,t}\cdot g_c(T_t,R_t,N_t),
\label{eq:q4_weather}
\end{equation}

其中 $T_t$、$R_t$、$N_t$ 为当日气温、降水与客流，$g_c(\cdot)$ 为品类弹性函数。
修正后异常天气日的补货偏差显著缩小。

\subsection{竞争对手定价}

现有方案 A 的加成率 $\eta_c$ 取历史均值，不随竞争环境调整。采集周边同类商超的
同品售价后，可定义竞争系数 $J_{c,t}$（本店价格与对手价格的相对位置），将加成率
修正为

\begin{equation}
\eta'_{c,t}=\eta_c\cdot(1-\lambda J_{c,t}),
\label{eq:q4_comp}
\end{equation}

其中 $\lambda\in(0,1)$ 为竞争响应强度。当对手降价（$J_{c,t}$ 增大）时本店加成率
相应下调，避免需求流失；对手提价时则可适度跟涨，实现差异化的动态定价。

\subsection{供货商数据}

现有模型假设任意数量即时到货。采集供货商的最小起订量 $MOQ_i$、交货提前期
$Ld_i$ 与历史准时率 $\tau_i$ 后，可加入可行性约束

\begin{equation}
x_i\ge MOQ_i\,y_i,
\label{eq:q4_moq}
\end{equation}

并以 $\tau_i$ 折减有效供给，按提前期滚动排产。修正后方案满足供货现实约束，避免
"模型给出 2.5 kg 起订、供应商要求 20 kg 起订"的不可执行情形。

\subsection{关联食材销售}

蔬菜常与肉类搭配烹饪，其需求受肉类消费带动。采集肉类销售流水后，计算肉类与
各蔬菜品类的相关系数 $R_c$，将需求曲线修正为

\begin{equation}
D'_c(P)=D_c(P)\left(1+R_c\,\frac{\Delta B_{\mathrm{meat}}}{B_{\mathrm{meat}}}\right),
\label{eq:q4_meat}
\end{equation}

其中 $\Delta B_{\mathrm{meat}}/B_{\mathrm{meat}}$ 为肉类销量的变动率。修正后定价
模型能捕捉荤素搭配的需求联动，节假日肉类促销期的蔬菜补货可提前加量。

综合以上六类数据，采集成本由低到高依次为供货商合同数据、库存盘点、气象接口、
客流计数、竞争对手调研与关联品类流水，商超可按投入产出比分批接入。每类数据都以
一个明确的参数或约束进入模型，修正路径清晰，为模型的持续改进预留了接口。
